\section{Discussion and Limitations}\label{sec:discussion}

\subsection{Fixed constitution versus adaptive institution}\label{sec:discussion:p2}

The architecture deliberately stops short of unrestricted
self-modification. Its constitution, transition topology, capability
matrix, and fixed verifier do not evolve at runtime. This limits the
system: a genuinely new role or transition requires a human amendment.
The same limit is also what makes every institutional change
interpretable. A candidate can improve within a known authority envelope,
but cannot redefine the envelope that judges it.

The constitution may still be wrong. Hash pinning makes amendment
explicit; it does not make the pinned rules wise. Human review, rationale,
and regression evidence therefore remain part of the trusted process.
Constitutional ARI-RQGM reduces the authority granted to model outputs; it
does not eliminate the responsibility of the people who define the law.

Emergency quarantine now preserves the epoch invariant. T16 immediately
force-closes the current epoch and atomically opens a new fingerprint while
quarantining the offending component. This makes the containment point
explicit in replay and prevents one fingerprint from naming two serving
sets. It does not solve adjudication denial of service or provide a
replacement role inside the emergency transaction.

\subsection{Instruction-following at authority seams}\label{sec:discussion:seams}

Most dangerous failures occur where a model-generated record crosses into
deterministic state. Schema validation alone is insufficient: a perfectly
formed record may be authored by the wrong role, request an illegal edge,
expand capability, refer to sealed text, or omit the provenance needed for
repair. The kernel therefore validates structure, author, timing,
authority, hash, and context as separate dimensions.

There is a complementary semantic limitation. A valid attack or judgment
may still be mistaken. Role separation and replay make the mistake
observable and contestable, but they do not turn model judgment into
ground truth. Fixed anchors are used where execution evidence or curated
labels exist; elsewhere, conclusions remain probabilistic and should not
be promoted to constitutional fact.

Likewise, ``research continues when governance fails'' is a conditional
availability policy, not a general safety principle. Continuation should
be limited to candidate generation and local computation in a sandbox.
Writes to external services, physical devices, or production data need a
separate fixed gate that fails closed on governance failure. The present
kernel does not enforce that environmental condition.

\subsection{Scaling limits}\label{sec:discussion:scaling}

Governance adds model calls, replay work, storage, and boundary latency.
Budgets can lower the governance level or suppress optional co-evolution,
but a heavily degraded run collects less evidence about its institution.
Large registries also increase the number of candidates, role interactions,
and dependency edges that repair must inspect.

The version-two audit digest binds the event type, identifier, transaction,
payload, and predecessor and therefore exposes replay-semantic interior
modification, but not rollback of
both the history and its local pin to the same old prefix. Strong rollback
resistance needs an externally stored signed head, periodic commitment to
an independent witness, or a monotonic counter. Permanent append-only
storage also conflicts with confidentiality, personal data, and deletion
requests. Sensitive plaintext should not enter the ledger without a
design such as key-destruction erasure or an auditable redaction record.

The present implementation assumes one authoritative event order per
checkpoint. Distributed executors would require a replicated log,
deterministic conflict resolution, and a clear definition of which epoch
was serving when concurrent artefacts were produced. The current
single-writer discipline is a sound local boundary, not yet a distributed
consensus protocol.

Data are another limit. Reviewer improvement needs held-out anchors;
clean-room regeneration needs sufficiently abstract failure summaries;
and utility-policy comparison needs a stable axis basis. When those inputs
are absent, the system degrades conservatively rather than inventing a
claim of superiority. Conservative degradation preserves integrity but
reduces the opportunity for adaptation.

Repeatedly selecting against the same validated-attack pool can overfit both
candidates and utility policy. A credible study must separate development,
selection, and secret final cases, account for adaptive repeated selection,
and compare same-model, cross-model, and cross-provider role assignments.
Record-level role separation does not imply independent errors when several
roles share a model family. Collusion, adjudication denial of service, and
biased anchors therefore remain empirical threats rather than solved
properties.

\subsection{Outlook}\label{sec:discussion:outlook}

The next empirical step is a pre-registered comparison across governance
configurations, task families, models, and seeds. The next systems step is
to make evidence coverage explicit: for each role, report which sanctions,
replay cases, anchors, and successor stages were actually exercised.
Longer runs are needed to study repeated policy rewrites and institutional
cycling.

The priorities are concrete: externally anchor audit heads and isolate role
and checkpoint authority at the process or operating-system boundary;
compare a fixed institution, original RQGM, and
the complete system with component ablations; measure attribution and
repair accuracy on labelled causal graphs; and quantify manuscript-gate
coverage together with governance time, token, and monetary cost.

The broader research question is whether a fixed constitution can support
open-ended institutional improvement without becoming either too weak to
constrain it or too rigid to permit it. Constitutional ARI-RQGM provides a
concrete, replayable platform on which that question can be measured. Its
present claim is intentionally narrower: model-defined research roles can
co-evolve while registry authority, amendment, provenance, and final
evidence verification remain outside their control.
